\chapter{PDE and Submartingale Characterizations}

\section{Reflected diffusions and parabolic PDEs}

Let $D\subset\R^d$ be a domain for which the reflected SDE is well posed in
law. In the time-homogeneous case, write
\begin{equation}
  X_t=x+\int_0^t a(X_s)\,\dd s
      +\int_0^t\sigma(X_s)\,\dd B_s
      +\int_0^t\gamma(X_s)\,\dd L_s, \label{eq:chapter3-rsde}
\end{equation}
where $\gamma$ is an inward reflection field and $\dd L$ is supported on the
boundary. Set
\[
  c(x):=\sigma(x)\sigma(x)^\top
\]
and define, for $\varphi\in C^2(\overline D)$,
\begin{equation}
  \mathcal A\varphi(x)
  :=a(x)\cdot\nabla\varphi(x)
    +\frac12\operatorname{tr}\!\left(c(x)D^2\varphi(x)\right). \label{eq:diffusion-generator}
\end{equation}
The oblique derivative at the boundary is
$\partial_\gamma\varphi:=\nabla\varphi\cdot\gamma$.

Uniqueness in law is needed here: weak existence alone does not select a
unique family of transition probabilities. Under well posedness, define
\[
  P_t\varphi(x):=\E_x[\varphi(X_t)].
\]

\begin{proposition}[A generator core candidate]
Assume $a,\sigma,$ and $\gamma$ are bounded and continuous. If
$\varphi\in C_b^2(\overline D)$ satisfies
$\partial_\gamma\varphi=0$ on $\partial D$, then
\begin{equation}
  \lim_{t\downarrow0}\frac{P_t\varphi(x)-\varphi(x)}{t}
  =\mathcal A\varphi(x),
  \qquad x\in\overline D. \label{eq:generator-at-zero}
\end{equation}
\end{proposition}

\begin{proof}
It\^o's formula gives
\begin{align*}
  \varphi(X_t)-\varphi(x)
  &=\int_0^t\mathcal A\varphi(X_s)\,\dd s
    +\int_0^t\nabla\varphi(X_s)\sigma(X_s)\,\dd B_s\\
  &\quad+\int_0^t\partial_\gamma\varphi(X_s)\,\dd L_s.
\end{align*}
The last term is zero because $\dd L$ is supported on $\partial D$, and the
stochastic integral is a true martingale because its integrand is bounded.
After taking expectations, continuity of $X$ and bounded continuity of
$\mathcal A\varphi$ give \eqref{eq:generator-at-zero} by dominated
convergence.
\end{proof}

\begin{remark}[What the proposition does not prove]
The proposition identifies a natural class of functions on which the
probabilistic generator acts as $\mathcal A$ with an oblique Neumann boundary
condition. It does not identify the whole generator domain, prove that this
class is a core, or prove strong continuity of $(P_t)$. A pointwise derivative
at $t=0$ cannot replace the separate Feller and strong-continuity arguments
required by semigroup theory.
\end{remark}

\subsection{The backward Kolmogorov equation}

For time-dependent coefficients, let $X^{t,x}$ be the reflected diffusion
started from $x$ at time $t$, and define
\begin{align*}
  \mathcal A_s\varphi(x)
  &:=a(s,x)\cdot\nabla\varphi(x)
    +\frac12\operatorname{tr}\!\left(
        c(s,x)D^2\varphi(x)
      \right),\\
  c(s,x)&:=\sigma(s,x)\sigma(s,x)^\top.
\end{align*}
The natural probabilistic object is now a transition family
$P_{t,s}$, not a one-parameter semigroup.

\begin{theorem}[Classical backward representation]
Let $u\in C^{1,2}([0,T]\times\overline D)$ satisfy
\begin{align*}
  \partial_tu(t,x)+\mathcal A_tu(t,x)&=0,
      &&(t,x)\in[0,T)\times D,\\
  \partial_{\gamma(t,\cdot)}u(t,x)&=0,
      &&(t,x)\in[0,T)\times\partial D,\\
  u(T,x)&=g(x), &&x\in\overline D.
\end{align*}
Assume that $u$, its derivatives appearing above, and the stochastic
integrand $\nabla u\,\sigma$ have enough growth control to justify It\^o's
formula and the martingale expectation. Then
\begin{equation}
  u(t,x)=\E\bigl[g(X_T^{t,x})\bigr]. \label{eq:backward-representation}
\end{equation}
\end{theorem}

\begin{proof}
Apply It\^o's formula to $u(s,X_s^{t,x})$ on $[t,T]$. The time--space drift
vanishes by the PDE, and the finite-variation boundary term vanishes because
$\partial_\gamma u=0$ where the regulator increases. The remaining
stochastic integral is a martingale. Taking expectations gives
\eqref{eq:backward-representation}.
\end{proof}

\begin{remark}
Existence of a classical solution and existence of a transition density are
separate analytic questions. On a sufficiently smooth domain, uniformly
parabolic operators with compatible H\"older data admit classical Neumann or
oblique-derivative theory. When a transition density exists, the function in
\eqref{eq:backward-representation} may be written as its integral against
$g$. A PDE solution is a deterministic function; it is not itself a Markov
process.
\end{remark}

\begin{theorem}[Feynman--Kac formula]
Let $r$ be bounded from below, and let
$u\in C^{1,2}([0,T]\times\overline D)$ solve
\begin{align*}
  \partial_tu+\mathcal A_tu-r u&=0
      &&\text{on }[0,T)\times D,\\
  \partial_\gamma u&=0
      &&\text{on }[0,T)\times\partial D,\\
  u(T,\cdot)&=g.
\end{align*}
Assume the exponential weight and the stochastic integral below are
integrable. Then
\begin{equation}
  u(t,x)
  =\E\!\left[
      g(X_T^{t,x})
      \exp\!\left(
        -\int_t^T r(s,X_s^{t,x})\,\dd s
      \right)
    \right]. \label{eq:feynman-kac-reflected}
\end{equation}
\end{theorem}

\begin{proof}
Apply the product rule to
\[
  \exp\!\left[
    -\int_t^s r(q,X_q^{t,x})\,\dd q
  \right]u(s,X_s^{t,x}).
\]
The PDE cancels the drift and the boundary condition cancels the regulator
term. The remaining local martingale is a true martingale under the stated
integrability assumption. Evaluating at $t$ and $T$ proves the formula.
\end{proof}

\subsection{Expected hitting times}

Let $K$ be an open set with $\overline K\subset D$, put
$U:=D\setminus\overline K$, and define
\[
  \tau_K:=\inf\{t\geq0:X_t\in\overline K\}.
\]

\begin{theorem}[Poisson problem for the mean hitting time]
Assume $D$ is bounded and that the stopped stochastic integrals below are
true martingales. Suppose $u$ extends to a $C^2$ function on a neighbourhood of
$\overline U$ and satisfies
\begin{align*}
  \mathcal Au&=-1 &&\text{in }U,\\
  \partial_\gamma u&=0 &&\text{on }\partial D,\\
  u&=0 &&\text{on }\partial K.
\end{align*}
Then $\tau_K<\infty$ almost surely and
\begin{equation}
  u(x)=\E_x[\tau_K],
  \qquad x\in\overline U. \label{eq:mean-hitting}
\end{equation}
\end{theorem}

\begin{proof}
It\^o's formula stopped at $\tau_K\wedge t$ yields
\[
  \E_x u(X_{\tau_K\wedge t})
  =u(x)-\E_x(\tau_K\wedge t).
\]
Because $u$ is bounded on $\overline U$, the right-hand side gives a uniform
bound on $\E_x(\tau_K\wedge t)$. Monotone convergence then implies
$\E_x\tau_K<\infty$, and hence $\tau_K<\infty$ almost surely. Continuity of
$X$ gives $X_{\tau_K}\in\partial K$; bounded convergence and the inner
Dirichlet condition now give \eqref{eq:mean-hitting}.
\end{proof}

\begin{example}[Reflected Brownian motion in an annulus]
Let $D=B(0,R)$ and $K=B(0,r)$ with $0<r<R$. Let $X$ be Brownian motion
normally reflected at $\partial D$, so its generator is
$\frac12\Delta$. The mean time to hit $\overline K$ is radial:
$u(x)=U(|x|)$. For $r\leq\rho\leq R$, it solves
\[
  U''(\rho)+\frac{d-1}{\rho}U'(\rho)=-2,
  \qquad U(r)=0,
  \qquad U'(R)=0.
\]
Thus, for $d\geq3$,
\begin{equation}
  u(x)
  =\frac{r^2-|x|^2}{d}
   +\frac{2R^d}{d(d-2)}
      \bigl(r^{2-d}-|x|^{2-d}\bigr), \label{eq:annulus-dge3}
\end{equation}
where the factor $R^d$ is required by the Neumann condition. For $d=2$,
\begin{equation}
  u(x)=\frac{r^2-|x|^2}{2}+R^2\log\frac{|x|}{r}. \label{eq:annulus-d2}
\end{equation}
For $d=1$, on either component of $(-R,R)\setminus[-r,r]$,
\[
  u(x)=r^2-|x|^2+2R(|x|-r).
\]
\end{example}

\section{The submartingale problem}

The submartingale formulation characterizes reflected diffusions without
choosing a Brownian motion in advance. Let
\[
  G=\{x\in\R^d:\phi(x)>0\},
  \qquad
  \partial G=\{x:\phi(x)=0\},
\]
where $\phi\in C_b^2(\R^d)$ and
$|\nabla\phi|$ is bounded away from zero on the boundary. Let
$a(t,x)$ and $c(t,x)$ be the interior drift and covariance, let
$\gamma(t,x)$ be a continuous inward field with
\[
  \inf_{t,x\in\partial G}
  \gamma(t,x)\cdot\nabla\phi(x)>0,
\]
and let $\rho(t,x)\geq0$ describe possible boundary holding. Define
\begin{align*}
  \mathcal A_t f(t,x)
  &:=a(t,x)\cdot\nabla f(t,x)
    +\frac12\operatorname{tr}
      \bigl(c(t,x)D^2f(t,x)\bigr),\\
  \mathcal J_t f(t,x)&:=\gamma(t,x)\cdot\nabla f(t,x).
\end{align*}

\begin{definition}[Submartingale problem]
A probability measure $P$ on
$C([t_0,\infty);\R^d)$ solves the submartingale problem for
$(a,c,\gamma,\rho)$ with initial point $x_0\in\overline G$ if, for the
coordinate process $X$,
\begin{enumerate}[label=(\roman*)]
  \item
  \[
    P\bigl(X_{t_0}=x_0,
      \ X_t\in\overline G\text{ for every }t\geq t_0\bigr)=1;
  \]
  \item for every compactly supported
  $f\in C^{1,2}([t_0,\infty)\times\R^d)$ satisfying
  \begin{equation}
    \rho(t,x)\partial_t f(t,x)+\mathcal J_t f(t,x)\geq0
    \quad\text{on }[t_0,\infty)\times\partial G, \label{eq:submart-boundary-test}
  \end{equation}
  the process
  \begin{align}
    f(t,X_t)-f(t_0,X_{t_0})
    -\int_{t_0}^t\1_{\{X_s\in G\}}
       \bigl(\partial_s f+\mathcal A_s f\bigr)(s,X_s)\,\dd s
       \label{eq:submart-process}
  \end{align}
  is a $P$-submartingale.
\end{enumerate}
\end{definition}

The event in (i) is a pathwise confinement statement. Requiring only
$P(X_t\in\overline G)=1$ separately for each fixed $t$ would be weaker in
form, even though continuity can sometimes be used to recover the pathwise
version from a countable dense set of times.

\begin{theorem}[Semimartingale representation]
Under the usual smoothness and nondegeneracy assumptions for the
submartingale problem, $P$ solves the problem above if and only if there is a
continuous adapted nondecreasing process $L$, with $L_{t_0}=0$ and
\begin{equation}
  \int_{t_0}^t\1_{\{X_s\in G\}}\,\dd L_s=0, \label{eq:l-support-ch3}
\end{equation}
such that
\begin{align}
  M_t
  &:=X_t-X_{t_0}
     -\int_{t_0}^t\1_{\{X_s\in G\}}a(s,X_s)\,\dd s
     -\int_{t_0}^t\gamma(s,X_s)\,\dd L_s \label{eq:submart-M}\\
  \langle M^i,M^j\rangle_t
  &=\int_{t_0}^t\1_{\{X_s\in G\}}c_{ij}(s,X_s)\,\dd s, \label{eq:submart-bracket}\\
  \int_{t_0}^t\1_{\{X_s\in\partial G\}}\,\dd s
  &=\int_{t_0}^t\rho(s,X_s)\,\dd L_s. \label{eq:elastic-relation}
\end{align}
Here $M$ is a continuous local martingale. Under boundedness, it is a true
martingale on finite intervals.
\end{theorem}

\begin{remark}
When $\rho\equiv0$, relation \eqref{eq:elastic-relation} says that the process
spends zero Lebesgue time on the boundary. Positive $\rho$ produces sticky or
elastic boundary behavior. The process $L$ is part of the representation and
must start at zero, increase only on the boundary, and use the bracket lower
limit $t_0$; omitting any of these conditions leaves the characterization
incomplete.
\end{remark}

\begin{theorem}[A safe well-posedness regime]
Suppose $G$ is a bounded $C^2$ domain. Assume that $a$ and a square root
$\sigma$ of $c$ are bounded and globally Lipschitz in $x$, continuous in
$t$, and that $c$ is uniformly elliptic. Assume $\gamma$ is bounded,
Lipschitz, and uniformly inward pointing.

For $\rho\equiv0$, the reflected SDE is pathwise unique and its law is the
unique solution of the submartingale problem. Its transition family is strong
Markov; equivalently, the time--space process is a homogeneous strong Markov
process.

Assume in addition that $a,c,\gamma,$ and $\rho$ are time homogeneous and
that $\rho$ is bounded, nonnegative, and Lipschitz. Let
$(\overline X,\overline L)$ denote the nonsticky reflected diffusion and set
\[
  A_s=s+\int_{t_0}^s\rho(\overline X_r)\,\dd\overline L_r,
  \qquad
  \tau_t:=\inf\{s\geq t_0:A_s>t\},
  \qquad t\geq t_0.
\]
Then $X_t:=\overline X_{\tau_t}$ yields the elastic process, and the
corresponding submartingale problem is well posed.
\end{theorem}

\begin{proof}[Proof outline]
The smooth-domain reflected SDE has a unique strong solution under the stated
Lipschitz assumptions. It\^o's formula proves that its law solves the
submartingale problem. Conversely, the semimartingale representation and the
martingale representation theorem recover a weak reflected SDE from any
solution of the submartingale problem; uniqueness in law identifies it. In the
time-homogeneous elastic case, the inverse clock $\tau$ slows physical time
exactly when $\overline X$ accumulates boundary regulator. The change-of-
variables formula then gives \eqref{eq:elastic-relation}. The construction is
reversible because $A$ is continuous and strictly increasing.
\end{proof}

\begin{example}[Why time-dependent test functions matter]
Take $G=(0,\infty)$, $a=0$, $c=1$, $\gamma=1$, and $\rho=0$.
Reflected Brownian motion solves the submartingale problem. Brownian motion
stopped upon first hitting zero does not solve it, because after the hitting
time a time-dependent test function can change through its boundary trace
while the interior integral in \eqref{eq:submart-process} is switched off.
Nevertheless, for every time-independent
$f\in C_c^2(\R)$ with $f'(0)\geq0$, the stopped process satisfies
\[
  f(X_t)-\frac12\int_0^t\1_{\{X_s>0\}}f''(X_s)\,\dd s
\]
as a submartingale (indeed, as a martingale). Therefore one cannot replace
the full boundary condition
$\rho\,\partial_t f+\mathcal Jf\geq0$ by a condition involving only
$f$ and its spatial derivatives.
\end{example}

\begin{remark}[Localization]
Boundedness assumptions may be localized. If two sets of coefficients agree
on $[t_0,t_0+\varepsilon]\times U$, where $U$ is a neighborhood of the
starting point, then their stopped submartingale problems agree up to the
first exit from that time--space region, provided the localized problem is
well posed.
\end{remark}
